% !TeX encoding = UTF-8
% !TeX spellcheck = en_US
%%%%%%%%%%%%%%%%%%%%%%%%%%%%%%%%%%%%%%%%%%%%%%%%%%%%%%%%%%%%%%%%%%%%%%%%%%%%%
%  enslyonreview-exemple.tex — FULLY WORKED EXAMPLE for the enslyonreview class
%
%  ⚠  This file is for REFERENCE only.
%     To write your own review, open and edit enslyonreview-gabarit.tex
%     instead. Copy-paste sections from here as needed.
%
%  Content reproduces, verbatim, "Antiviral restriction factors in immunity
%  and disease" by Erika Larrahondo Rodríguez (M2 Review, Dec 2024) — the
%  finished-document example supplied by F. Roudier alongside review_M2.ott —
%  so the class can be checked page-for-page against a real, previously
%  approved review. Figures 1 and 2 are the original images extracted from
%  that PDF; Table 1 and the 50-entry reference list are retyped in full,
%  and citations are REAL biblatex \cite calls against
%  enslyonreview-exemple.bib (clickable, auto-numbered) rather than typed
%  bracket numbers — see that .bib file for the field conventions to follow.
%
%  Project structure:
%    enslyonreview-gabarit.tex ← YOUR blank template (start here)
%    enslyonreview-exemple.tex ← this file (full worked example)
%    enslyonreview.cls         ← class file (same directory or on TEXMFHOME)
%    enslyonreview.dbx         ← extra biblatex fields (pmid/pmcid/interest)
%    enslyonreview-exemple.bib ← the 50 references cited below
%    figures/                  ← fig1-restriction-mechanisms.jpg,
%                                 fig2-disease-drivers.jpg
%
%  Compile order:
%    pdflatex enslyonreview-exemple → biber enslyonreview-exemple
%    → pdflatex enslyonreview-exemple → pdflatex enslyonreview-exemple
%  On Overleaf: compilation is automatic (biber is used by default).
%%%%%%%%%%%%%%%%%%%%%%%%%%%%%%%%%%%%%%%%%%%%%%%%%%%%%%%%%%%%%%%%%%%%%%%%%%%%%

% ── THEME — orange (default, official review_M2.ott rust/peach palette) ────
\documentclass{enslyonreview}

%%%%%%%%%%%%%%%%%%%%%%%%%%%%%%%%%%%%%%%%%%%%%%%%%%%%%%%%%%%%%%%%%%%%%%%%%%%%%
%  METADATA — populates the full-width title block on page 1, the running
%  header/footer on every page, and the PDF metadata.
%%%%%%%%%%%%%%%%%%%%%%%%%%%%%%%%%%%%%%%%%%%%%%%%%%%%%%%%%%%%%%%%%%%%%%%%%%%%%

% ── TITLE — full title (page 1) + short form (running header) ─────────────
\reviewtitle{Antiviral restriction factors in immunity and disease: Regulation and dysregulation of their effector functions}
\reviewshorttitle{Antiviral Restriction Factors}

% ── AUTHOR — full name (page 1) + surname (running header) ────────────────
% \textsc{} on the surname (enslyonstage convention) — no effect on
% \reviewauthorshort, which stays plain (already just the surname, shown
% small in the italic running header).
\reviewauthor{Erika \textsc{Larrahondo Rodr\'iguez}}
\reviewauthorshort{Larrahondo Rodr\'iguez}

% ── AFFILIATION AND DATE — shown together in the peach band under the title
\reviewaffiliation{Master Biosciences, D\'epartement de Biologie, \'Ecole Normale Sup\'erieure de Lyon.}
\reviewdate{Dec 2024}

% ── KEYWORDS ────────────────────────────────────────────────────────────
\reviewkeywords{Antiviral, disease, innate immunity, restriction factors}

% ── ABSTRACT — single paragraph, no bibliographic reference ───────────────
\reviewabstract{Antiviral factors from the cell autonomous innate immunity are essential
  cellular proteins that inhibit viral replication and play a central role in the host's
  defense against infections. Novel technologies now enable the high-throughput screening of
  these antiviral factors and facilitate detailed characterization of their molecular
  mechanisms. In this review, we explore the diverse effector functions of these proteins,
  focusing on their restriction of molecular processes and their interactions with viral
  components. We also highlight how the dysregulation or genetic mutations of antiviral
  factors can contribute to the onset and development of diseases. Finally, we emphasize
  their dual role in immunity and pathology and discuss how interdisciplinary approaches will
  continue to advance the field.}

% ── BIBLIOGRAPHY — biblatex + biber (see compile order above) ─────────────
\addbibresource{enslyonreview-exemple.bib}

% ── "OF SPECIAL / OUTSTANDING INTEREST" ANNOTATIONS ────────────────────────
% Declared here rather than in the .bib — see enslyonreview-exemple.bib's
% header comment for why. Each key below matches an entrykey in that .bib
% exactly (case-sensitive); \refspecialinterest → one bullet (•),
% \refoutstandinginterest → two bullets (••), printed under that entry in
% the reference list.
\refspecialinterest{wangGenomewideCRISPRScreenings2024}{Using a genome-wide CRISPR-Cas9 screening, the authors identify SMCHD1 as a restriction factor for adeno-associated virus (AAV). Then they perform mechanistic studies that demonstrate SMCHD1 represses AAV transcription. This study identifies a restriction factor using a high-throughput screening and details its molecular mechanism.}

\refoutstandinginterest{luIFI16PromotesProgression2024}{Using in vitro and in vivo experiments, the authors demonstrate that IFI16 promotes clear cell renal cell carcinoma (ccRCC) cell proliferation and tumor growth. Using transcriptome analysis and experimental validation they confirm the effect relies on the PI3K/AKT pathway. This study details the molecular mechanism through which IFI16 acts as a cancer driver in ccRCC.}

\refspecialinterest{maharanaSAMHD1ControlsInnate2022}{Using in vitro experiments, authors show that SAMHD1 functions as a single-stranded RNA (ssRNA) 3' exonuclease and that in the absence of SAMHD1, cellular RNA is accumulated and leads to immune activation. This study demonstrates the role of SAMHD1 as a regulator of cellular RNA homeostasis.}

\refoutstandinginterest{naumannDNADeaminationRequired2023}{Using in vivo experiments, the authors demonstrate that APOBEC3A can drive tumor development and that this effect is compromised in the absence of its DNA deamination activity. This study demonstrates APOBEC3A drives tumor development in vivo through the same mechanism as antiviral restriction.}

\refoutstandinginterest{wuPhosphoproteomicsRevealsRole2022}{By performing phosphoproteomics on cells of chronic lymphocytic leukemia (CLL) patients and further experiments on KAP1 mutants, the authors show that phosphorylation of KAP1 though the B cell receptor pathway accelerates the cell cycle in CLL cells. This study shows a link between immune signaling, a restriction factor and cancer progression.}

\refspecialinterest{yangKAP1StabilizesMYCN2024}{Using in vitro and RNA stability experiments, authors find that KAP1 is upregulated by MYCN and by forming a complex with other proteins, stabilizes MYCN mRNA and therefore promotes neuroblastoma proliferation. This study illustrates a KAP1 cancer-driving activity dependent on a mechanism different to the one involved in viral restriction.}

\refspecialinterest{legrandSAMD9LActsAntiviral2024}{Here the authors use in vitro experiments and bioinformatic analyses to show that SAMD9L restricts primate lentiviruses by inhibiting translation through a conserved SLFN-like box active site. This study gives novel insight into the molecular mechanism through which SAMD9L restricts the viral cycle.}

\refspecialinterest{manjunathAPOBEC3BDrivesPKRmediated2023}{In this study the authors use different RNA viruses to show that APOBEC3B promotes the PKR signaling pathway and limits stress granule dissociation during viral infection. Here authors study how the interaction between two restriction factors participates in antiviral response and specifically translation inhibition.}

\refspecialinterest{okamotoSLFN11PromotesStalled2021a}{Using genome editing on a cell line derived from a Fanconi anemia (FA) patient, authors demonstrate that SLFN11 depletion lowers interstrand crosslink-induced chromosomal breaks and that SLFN11 promotes stalled fork degradation caused by nucleases in FA. This study demonstrates a disease-driving activity for SLFN11.}

\refspecialinterest{heBST2InducedMacrophage2023}{Through in vivo and in vitro experiments, the authors show that BST2 silencing restricts colorectal cancer progression and macrophage polarization to an M2 phenotype. Also, that depletion of tumor-associated macrophages alleviates the protumoral effect of BST2. This study details the mechanism by which BST2 drives colorectal cancer progression.}

\refspecialinterest{leonhardtAntiviralHIV1SERINC2023}{In this study the authors use cryoEM and fluorescent proteoliposome assay to determine that membrane transporters that can induce the flipping of lipids and therefore disrupt membrane asymmetry. This study uses structural and biochemical assays to unravel the viral restriction mechanism of SERINC proteins.}


%%%%%%%%%%%%%%%%%%%%%%%%%%%%%%%%%%%%%%%%%%%%%%%%%%%%%%%%%%%%%%%%%%%%%%%%%%%%%
\begin{document}
%%%%%%%%%%%%%%%%%%%%%%%%%%%%%%%%%%%%%%%%%%%%%%%%%%%%%%%%%%%%%%%%%%%%%%%%%%%%%

% ── TITLE BLOCK — must be the first call after \begin{document}; assembles
% the metadata above into the full-width header, then switches to two
% columns for everything that follows.
\maketitle

%%%%%%%%%%%%%%%%%%%%%%%%%%%%%%%%%%%%%%%%%%%%%%%%%%%%%%%%%%%
\section{Introduction}
%%%%%%%%%%%%%%%%%%%%%%%%%%%%%%%%%%%%%%%%%%%%%%%%%%%%%%%%%%%

Early and effective response against infection is crucial in antiviral immunity. Antiviral
factors, also referred to as restriction factors, are cellular proteins that confer intrinsic
or cell-autonomous innate immunity to viral pathogens by inhibiting viral entry or
replication. Restriction factors are a group of proteins that are able to restrict the viral
replication cycle through a wide variety of mechanisms, ranging from transcriptional
silencing to degradation of cellular biomolecules. Here we focused on restriction factors
whose activity has been described in mammals. Although the precise definition of restriction
factors has been debated, they share a set of common features: they have direct antiviral
activities and generally not part of a complex signaling cascade, they can be constitutively
expressed in most cell types and are often upregulated by interferon (IFN) (i.e a secreted
cytokine with antiviral functions), and they often present signatures of positive selection
(i.e the rapid fixation of beneficial non-synonymous mutations) as a result of the arms race
between host and virus \cite{blanco-meloIntrinsicCellularDefenses2012,klugeSnapShotAntiviralRestriction2015}.

The term ``restriction factor'' was first introduced in the early 1970s when the Friend virus
susceptibility protein 1 (Fv1) was identified as restricting the retroviral Murine Leukemia
Virus (MLV) infection in mice \cite{lillyFv2IdentificationLocation1970}. Since then, cell-intrinsic restriction factors
have been identified to interfere with virtually all steps of the viral cycle and all virus
types as part of the first line of defense against these pathogens. Several of these factors,
such as SAMHD1 \cite{cogginsSAMHD1FunctionsHuman2020} and Tetherin/BST-2 \cite{sauterCounteractionMultifunctionalRestriction2014}, have been extensively
studied, with their molecular mechanisms and corresponding viral resistance strategies
described in detail. However, new restriction factors are still actively being discovered and
technologies such as genome-wide CRISPR screenings have facilitated the identification of
these proteins \cite{wangGenomewideCRISPRScreenings2024}, which will surely lead to the discovery of unexpected
restriction mechanisms.

Although restriction factors have often been classified according to the step of the viral
cycle they inhibit \cite{klugeSnapShotAntiviralRestriction2015,serreroRestrictionFactorsRegulating2024}, many of these proteins exert their inhibitory
action through common effector functions, such as nucleic acid degradation or
protein-protein interactions. For instance, RNA-degrading restriction factors may act either
by directly degrading viral RNA or by targeting viral mRNA, thereby preventing the synthesis
of viral proteins. Moreover, these restriction factors are often broadly effective, inhibiting
not only a wide array of viruses but also diverse viral families \cite{klugeSnapShotAntiviralRestriction2015}. However,
these powerful effector functions can also be directed against endogenous cellular components
or processes under physiological co
nditions, and may therefore be harmful to the host.
Highlighting the molecular mechanisms by which restriction factors antagonize the viral cycle
can enhance our understanding of how these proteins successfully inhibit different viruses and
how the dysregulation of innate restriction factors can have deleterious effects on the host.

As studies of antiviral responses in different host-pathogen models progress, new restriction
factors and associated mechanisms are characterized. Moreover, new insights arise on how
restriction factors can be implicated in diseases, notably in inflammatory diseases, cancer
and autoimmune disorders \cite{luIFI16PromotesProgression2024,russellSAMD9LAutoinflammatoryAtaxia2021,morelPKRDoubleStranded2009}. In this review, we explore
different effector mechanisms through which restriction factors inhibit viral replication
cycles and describe dysregulations in their effector functions and the resulting detrimental
effects in the host. We aim to provide specific examples of restriction factors acting as
drivers of disease and illustrate how these disease-driving mechanisms could be related or not
with their antiviral activity. Finally, we offer our perspective on how restriction factors
can both inhibit the viral cycle and act as drivers of disease, and how the field continuously
evolves towards a more integral view of restriction factors.

% A "key definitions" box for conceptual terms used throughout the review
% — a single-column reviewbox (not reviewbox*) placed right after the
% Introduction, as is common in Current Opinion-style reviews. The
% definitions themselves are not new: each restates, as a standalone
% glossary entry, what the Introduction paragraph above already says
% parenthetically inline.
\begin{reviewbox}[!ht]
  \caption{\textbf{Key definitions.}}
  \label{box:definitions}
  \begin{description}
    \item[Restriction factor] Cellular protein that confers intrinsic,
      cell-autonomous innate immunity by directly inhibiting viral entry
      or replication.
    \item[Effector mechanism] The specific molecular process (nucleic
      acid degradation, transcription/translation inhibition, direct
      protein--virus interaction\ldots) through which a restriction
      factor exerts its antiviral activity.
    \item[Interferon (IFN)] Secreted cytokine with antiviral functions;
      many restriction factors are constitutively expressed and further
      upregulated by IFN signaling.
    \item[Positive selection] The rapid fixation of beneficial
      non-synonymous mutations -- a genomic signature of the evolutionary
      arms race between host restriction factors and viruses.
  \end{description}
\end{reviewbox}

%%%%%%%%%%%%%%%%%%%%%%%%%%%%%%%%%%%%%%%%%%%%%%%%%%%%%%%%%%%
\section{Nucleotide or nucleic acid degradation}
%%%%%%%%%%%%%%%%%%%%%%%%%%%%%%%%%%%%%%%%%%%%%%%%%%%%%%%%%%%
% Each effector-mechanism section below is a self-contained development —
% no \label needed unless another section explicitly cross-references it
% (\cref{sec:...} / \crefcolor{sec:...}).

An important mechanism that inhibits the viral cycle is the degradation of viral RNA/DNA or
the prevention of viruses from synthesizing new genetic material (\crefcolor{fig:mechanisms}).
A well studied example of this mechanism is the Sterile Alpha Motif and HD domain-containing
protein 1 (SAMHD1), a deoxynucleoside triphosphate triphosphohydrolase that hydrolyzes dNTPs
into deoxyribonucleosides and inorganic triphosphates. This protein is widely expressed at a
basal level and highly expressed in myeloid lineage cells. It is also known to restrict HIV-1
by inhibiting reverse transcription and viral complementary DNA (cDNA) synthesis through the
depletion of cellular dNTPs \cite{goldstoneHIV1RestrictionFactor2011}. After this catabolic activity had been
described, SAMHD1 was also shown to act as a restriction factor via RNase activity by
degrading genomic RNA of HIV-1 and other retroviruses \cite{choiSAMHD1SpecificallyRestricts2015}. However, the role of
RNA degradation in SAMHD1-mediated restriction has been debated, and a recent study has shown
that RNA-bound SAMHD1 does not have a dNTPase or RNase activity \cite{orrisGuaninecontainingSsDNARNA2023}. The precise
role of the catalytic activities of SAMHD1 in viral restriction remains to be elucidated.

Cellular-intrinsic proteins that show antiviral activity through the targeting of nucleic
acid material can also be drivers of disease. For instance, genetic mutations in SAMHD1 cause
an inflammatory encephalopathy called Aicardi--Goutières syndrome (AGS). The specific
consequences of defects in SAMHD1 function in this disease are still being studied. Recently,
SAMHD1 has been shown to regulate innate immunity by lowering cellular RNA levels and
therefore inhibiting the release of immunogenic double-stranded RNA from RNA-protein
condensates. The mutations in SAMHD1 could therefore lead to an activation of immune response
against cellular RNA and be linked to the inflammation and neurodegeneration seen in AGS
\cite{maharanaSAMHD1ControlsInnate2022}. Moreover, restriction factors are sometimes associated with cancer
progression, often because of their ability to disrupt physiological processes. Multiple
studies have aimed to determine the role of SAMHD1 in cancer, with different roles being
observed depending on the cancer type. In ovarian cancer, SAMHD1 depletion seems to induce
upregulation of proinflammatory cytokines and an activation of innate immune cell signaling,
and is consequently linked with improved prognosis \cite{gutierrez-chamorroSAMHD1ExpressionModulates2023}. In solid
cancers, a positive prognosis has also been associated with low or no expression of SAMHD1,
explained by its ability to modulate DNA damage response and a consequent increase in
$\gamma$-H2AX and apoptosis in SAMHD1-depleted cells (\crefcolor{fig:disease}A)~\cite{felipModulationDNADamage2022}.

Other examples of restriction factors that inhibit the viral cycle through modification or
degradation of nucleic acids include RNaseL, Zinc-finger Antiviral Protein (ZAP) and the
Apolipoprotein B mRNA Editing Catalytic Polypeptide-like (APOBEC) family. RNaseL is able to
cleave viral RNA when activated by the 2$'$-5$'$-linked oligoadenylates produced by the
2$'$-5$'$-oligoadenylate synthetase 1 (OAS1) protein \cite{klugeSnapShotAntiviralRestriction2015}. Recently, OAS1
gain-of-function variants have been shown to cause an autoinflammatory immunodeficiency
through RNAse L-mediated dysfunctions, that include cleavage of cellular RNA. This
immunodeficiency can be cured by inhibiting RNAseL \cite{maggHeterozygousOAS1Gainoffunction2021}. ZAP, also known as
poly(ADP-ribose) polymerase-13 (PARP13) selectively binds and degrades viral RNA, and is
active against a wide variety of viruses, including HIV-1 and hepatitis B virus
\cite{shaoVersatilityZincFingerAntiviral2024}. Finally, members of the APOBEC family induce the deamination of cytosines to
uracils and consequently introduce point mutations in viral genomes. The full scope of their
antiviral activity still remains to be discovered, with new antagonizing effects against
viruses still being reported \cite{gottigApobec3ADeaminationFunctions2023}. In this family, APOBEC3A drives tumor
development by a DNA deamination-dependent mechanism \cite{naumannDNADeaminationRequired2023}.

%%%%%%%%%%%%%%%%%%%%%%%%%%%%%%%%%%%%%%%%%%%%%%%%%%%%%%%%%%%
\section{Transcription inhibition}
%%%%%%%%%%%%%%%%%%%%%%%%%%%%%%%%%%%%%%%%%%%%%%%%%%%%%%%%%%%

In the presence of viral DNA, another restriction mechanism is the inhibition of viral gene
transcription (\crefcolor{fig:mechanisms}). The KRAB-associated protein 1 (KAP1), also known as
tripartite motif containing 28 (TRIM28) is a core component of the constitutive
heterochromatin machinery and recruits repressive histone modifiers that silence genes
through H3K9 trimethylation. In addition, KAP1-mediated DNA methylation can also silence the
transcription of viral genes in both RNA and DNA viruses. It was recently demonstrated that
in Epstein-Barr Virus (EBV), KAP1 interacts with the $\gamma$-interferon-inducible protein 16
(IFI16) to silence the EBV lytic switch protein ZEBRA, facilitating its heterochromatinization
and leading to viral latency \cite{xuIFI16PartnersKAP12022}. For Moloney murine leukemia virus (MoLV)
silencing, KAP1 is recruited to the provirus 5$'$ long-terminal, which results in a potent
transcriptional silencing. The latter is dependent on the K779 residue of the protein
\cite{leeCharacterizationInteractionTrim282018}. The extent of KAP1 antiviral activity remains to be elucidated in other
viruses, notably in HIV where its role remains elusive.

KAP1 also illustrates how dysregulation of transcription-inhibiting restriction factors can
lead to detrimental phenotypes. A recent pan-cancer analysis showed that KAP1 is more
expressed in various tumor tissues than in normal tissues, and that its abnormal expression is
associated with poorer prognosis. Specifically, higher expression of KAP1 was positively
correlated with tumor mutational burden and microsatellite instability \cite{shangSystematicPancancerAnalysis2023}. KAP1
parental and maternal pathogenic variants have also been reported to cause renal embryonal
tumors, probably through a deregulated transcriptional landscape during nephrogenesis that
ultimately leads to tumorigenesis \cite{whitworthWilmsTumourResulting2024}. Additionally, constitutive KAP1
phosphorylation through the B-cell receptor signaling pathway was higher in chronic
lymphocytic leukemia (CLL) patients compared to control patients, with KAP1 depletion
decelerating cell cycle progression of CLL cells (\crefcolor{fig:disease}B) \cite{wuPhosphoproteomicsRevealsRole2022}.
Importantly, the mechanism by which restriction factors inhibit the viral cycle can be
different to the mechanism by which they are involved in diseases. For example, KAP1 was
recently reported to stabilize the mRNA of the MYCN proto-oncogene, with KAP1 knockdown
decreasing MYCN-mediated neuroblastoma progression \cite{yangKAP1StabilizesMYCN2024}.

MxA and MxB proteins are large GTPases of the dynamin superfamily that also act as restriction
factors by inhibiting viral transcription. MxA acts as a restriction factor against a wide
range of viruses including orthomyxoviruses and paramyxoviruses. Although known to accumulate
in the cytoplasm, MxA was also shown to mediate transcription inhibition of the influenza
virus genome by forming a complex with the viral nucleoprotein (NP) \cite{turanNuclearMxAProteins2004}. Recent
studies on the link between MxA and disease have mostly focused on its influence on
sensitivity to chemotherapeutic agents or its relationship with prognosis
\cite{hayesExpressionMxAEsophageal2024,liExpressionMxAProtein2020}. MxB, another large GTPase, restricts various viruses through multiple
antiviral mechanisms, including significantly reducing hepatitis B virus and herpesvirus RNA
transcription, as well as inhibiting proviral DNA integration
\cite{xieMxBImpedesNUP358mediated2020,wangInterferoninducibleMX2Host2020,schillingHumanMxBProtein2018}. Besides MxA binding of a viral protein and MxB
restriction that seems to be dependent on GTPase activity, other mechanisms of transcription
inhibition have also been described. In addition to its role in EBV latency, IFI16 is a
nuclear DNA sensor that binds the herpes simplex virus 1 (HSV-1) genome in a
sequence-independent manner \cite{howardNuclearDNASensor2022}. IFI16 accumulates on the genome by interacting
with the backbone of dsDNA, reduces the association of transcription factors and leads to
gene repression \cite{johnsonIFI16RestrictsHSV12014}. A wide set of studies on the role of IFI16 in driving
cancers, mostly through immune signaling, have been performed, and showed recently that it
promotes clear renal cell carcinoma due to increased levels of IL6 \cite{luIFI16PromotesProgression2024}.

% ── FULL-WIDTH FIGURE — figure* spans both columns; a plain figure (no
% star) would stay within a single column instead. \label first, referenced
% from the text above/below via \crefcolor{fig:mechanisms}.
\begin{figure*}[!t]
  \centering
  \includegraphics[width=0.92\textwidth]{fig1-restriction-mechanisms}
  \caption{\textbf{Restriction factors inhibit the viral cycle through different effector
    mechanisms.} Two RNA viruses (green: retrovirus / yellow: ssRNA virus) and one DNA virus
    (white) are used to illustrate restriction factor activity through the different stages
    of the viral cycle. Blue cells represent host cells. Only nucleic acids from viral origin
    are represented. Restriction factors are depicted in red. Arrows indicate binding, dashed
    arrows indicate indirect effects.}
  \label{fig:mechanisms}
\end{figure*}

%%%%%%%%%%%%%%%%%%%%%%%%%%%%%%%%%%%%%%%%%%%%%%%%%%%%%%%%%%%
\section{Translation inhibition}
%%%%%%%%%%%%%%%%%%%%%%%%%%%%%%%%%%%%%%%%%%%%%%%%%%%%%%%%%%%

Restriction factors can also inhibit the viral cycle by affecting viral protein translation
(\crefcolor{fig:mechanisms}). The sterile alpha motif domain-containing 9 and 9-like proteins
(SAMD9/9L) are encoded by paralogous genes and are ubiquitously expressed in many human
tissues, where they are known to be poxvirus restriction factors that inhibit translation
\cite{mengParalogousPairMammalian2018}. SAMD9 induces stress antiviral granules that inhibit translation by a
mechanism including inhibition of ribosome entry and sequestration of viral RNA
\cite{sivanHumanHostRange2018}. Recently, SAMD9L but not SAMD9 has also been shown to restrict HIV-1 and
primate lentiviruses through the inhibition of viral and cellular translation. SAMD9L inhibits
the late stages of viral replication in interferon-stimulated cells through restriction of
viral protein synthesis. Interestingly, this inhibition was mediated by a Schlafen-like active
site, identified by homology with the SLFN11, SLFN12, and SLFN13 Schlafen proteins that also
act as cell autonomous restriction factors by inhibiting translation \cite{legrandSAMD9LActsAntiviral2024}.

Pathogenic variants of SAMD9 are the cause of a complex multisystem disorder named MIRAGE
syndrome (myelodysplasia, infection, restriction of growth, adrenal hypoplasia, genital
phenotypes, and enteropathy), which was first described in 2016 and is due to gain-of-function
mutations in the SAMD9 growth repressor gene that generate tissue hypoplasia. This syndrome is
still being studied for the large range of phenotypic manifestations, as well as the different
missense variants that drive the phenotypes \cite{suntharalinghamEmergingPhenotypesLinked2022}. SAMD9L mutations can
cause severe autoinflammatory disease and ataxia-pancytopenia. Although SAMD9L already
inhibits translation in physiological conditions, this activity is greatly increased both by
missense and by truncating mutations. Moreover, phenotypic alterations triggered by truncating
mutations might indicate an allosteric regulation of SAMD9L by its C terminus that, when lost,
can lead to the exaggerated translational repression (\crefcolor{fig:disease}C) \cite{russellSAMD9LAutoinflammatoryAtaxia2021}.
The dramatic phenotypes caused by SAMD9/9L mutations highlight the importance of some
restriction factors not only in innate immunity proteins but also as regulators of
physiological processes that upon deregulation can lead to deleterious phenotypes.

Numerous other proteins are both able to inhibit translation of viral proteins and implicated
in disease. Protein kinase R (PKR) is an interferon-induced kinase activated by binding to
viral dsRNA and known to impair viral translation through the phosphorylation of the
translation initiation factor eIF2$\alpha$. Interestingly, this translational shutdown has
been shown to be driven by APOBEC3B \cite{manjunathAPOBEC3BDrivesPKRmediated2023}. The activation of PKR has been
linked to neurodegenerative disorders and has been hypothesized to be an initiator of
deregulated translation in Alzheimer's disease lymphocytes via p53 \cite{morelPKRDoubleStranded2009}. As
mentioned above, members of the Schlafen family of proteins can also downregulate viral
protein synthesis. For example, SLFN11 has been proven to restrict flavivirus and HIV-1
infection, in a mechanism that seems to involve tRNA nucleolytic activity, as shown for
SLFN13 \cite{valdezSchlafen11Restricts2019}. SLFN11 is also a putative DNA/RNA helicase considered as a guardian
of the genome that has been shown to destabilize stalled replication forks and therefore
drive Fanconi anemia \cite{okamotoSLFN11PromotesStalled2021a}. The mechanisms by which the wide range of translation
inhibitory proteins can trigger diseases still remain largely unknown.

%%%%%%%%%%%%%%%%%%%%%%%%%%%%%%%%%%%%%%%%%%%%%%%%%%%%%%%%%%%
\section{Direct protein-virus interaction}
%%%%%%%%%%%%%%%%%%%%%%%%%%%%%%%%%%%%%%%%%%%%%%%%%%%%%%%%%%%

The best characterized example of a protein inhibiting the viral cycle only through physical
interaction with viral components is perhaps tetherin (\crefcolor{fig:mechanisms}). Also known as
bone marrow stromal antigen 2 (BST2) or CD317, tetherin was named after its ability to
physically tether virions to the plasma membrane of infected cells and therefore prevent viral
particle release. Tetherin is a cell membrane embedded dimer composed of a cytoplasmic region,
a transmembrane helix, an extracellular coiled coil domain, and a C-terminal membrane anchor.
The C-terminal anchor of tetherin is responsible for the attachment to the budding virus and
it is the ectodomain that bends and allows the transition between the membrane bound tetherin
and the bridge structure between the membrane and the virion \cite{ozcanBendingBST2Coiledcoil2017}. The antiviral
activity of tetherin has been documented for an extensive range of viruses that include HIV-1,
influenza A virus, and vesicular stomatitis virus. This list of targeted viruses will continue
to extend, as exemplified by recent work proving rabies virus inhibition via this mechanism
\cite{tanwattanaHumanBST2Inhibits2023}.

In the past years, a growing body of evidence showed that tetherin can have an oncogenic role.
Importantly, as a transmembrane glycoprotein, tetherin is also expressed in differentiated
human B cells and enhances cell-cell interaction in the pre-B cell stages of differentiation
\cite{jinBST2PromotesGrowth2019}. In colorectal cancer, in vivo and in vitro experiments showed that upregulation
of tetherin increased the infiltration of tumor-associated macrophages, polarizing them to an
M2 phenotype which in turn leads to an immunosuppressive tumor microenvironment
\cite{heBST2InducedMacrophage2023}. In oral squamous cell carcinoma (OSCC), tetherin was shown to activate the
epidermal growth factor receptor (EGFR) signaling pathway, important for proliferation and
apoptosis in mammalian cells. Specifically, tetherin overexpression enhanced OSCC cell
proliferation, inhibited OSCC cell apoptosis and induced resistance to the gefitinib
antitumor agent (\crefcolor{fig:disease}D) \cite{jinBST2PromotesGrowth2019}.

% Single-column-wide figure content on a figure* (two-column) float — the
% *environment* controls column span (both here), \includegraphics's own
% width= controls the image's actual size within it.
\begin{figure*}[!t]
  \centering
  \includegraphics[width=0.55\textwidth]{fig2-disease-drivers}
  \caption{\textbf{Restriction factors as drivers of disease.} The potential disease-driving
    roles of one protein from each type of effector mechanism are summarized. \textbf{A.}
    SAMHD1. \textbf{B.} KAP1. \textbf{C.} SAMD9L \textbf{D.} Tetherin. Blue cells represent
    host cells. Red cells represent cancer cells. Red shading symbolizes dysfunction or
    dysregulation. In immune cells, the red outline indicates activation. OSCC: Oral squamous
    cell carcinoma. CLL: chronic lymphocytic leukemia. AGS: Aicardi--Goutières syndrome.}
  \label{fig:disease}
\end{figure*}

Additional restriction factors are also able to bind viral components to restrict the viral
cycle. Two members of the interferon-induced transmembrane (IFITM) family of proteins, IFITM2
and IFITM3, have been shown to antagonize the HIV-1 envelope glycoprotein (Env) by interacting
with it in cells and thereby impairing correct virion assembly \cite{yuIFITMProteinsRestrict2015}. IFITM3 has been
identified as a $\gamma$-secretase modulatory protein in Alzheimer's disease, with experiments
showing that knockout of IFITM3 reduces $\gamma$-secretase activity and thereby the formation
of amyloid plaques. These results point to IFITM3 as a link between neuroinflammation and
Alzheimer's disease \cite{hurInnateImmunityProtein2020}. The serine incorporator SERINC3 and SERINC5 proteins have
recently been reported to restrict viral infection after the presence in the envelopes of
HIV-1 virions was shown to reduce infectivity and to disrupt virus membrane asymmetry
\cite{leonhardtAntiviralHIV1SERINC2023}. Interestingly, SERINC proteins have not yet been reported as drivers of
disease. The P-selectin glycoprotein ligand 1 (PSGL-1) protein also restricts HIV-1 infection
by binding and sequestering the gp41 protein at the plasma membrane, inhibiting envelope
incorporation in the virions \cite{liuPSGL1InhibitsHIV12020}. The PSGL-1 protein has been proposed as an immune
checkpoint and in an endotoxemia mice model, the blockage of PSGL-1 by an antibody
significantly increased survival and decreased lung injury and levels of inflammatory factors
\cite{wangRolePSGL1Pathogenesis2019}.

% ── FULL-WIDTH TABLE — table* (two columns); table alone would stay within
% one column, exactly like figure/figure* above. \footnotesize and
% \renewcommand{\arraystretch}{1.3} are plain LaTeX, not class macros: a
% four-column table this dense needs a smaller font and more row padding
% than the default to stay readable — tune both to taste for your own table.
% The >{\justifying\parindent0pt\arraybackslash} column prefixes (ragged2e,
% loaded by the class) give properly justified, hyphenated cell text
% without a paragraph indent — the plain p{width}/X column default is
% justified too, but WITH a first-line indent inherited from the
% document's \parindent, which looks wrong inside a table cell.
\begin{table*}[!t]
  \caption{\textbf{Antiviral restriction factors and their roles as disease drivers.}}
  \label{tab:summary}
  \footnotesize
  \renewcommand{\arraystretch}{1.3}
  \begin{tabularx}{\textwidth}{|>{\justifying\parindent0pt\arraybackslash}p{2.6cm}|>{\justifying\parindent0pt\arraybackslash}p{2.7cm}|>{\justifying\parindent0pt\arraybackslash}p{3.6cm}|>{\justifying\parindent0pt\arraybackslash}X|}
    \hline
    \rowcolor{colorAbstractBg}
    \textbf{Type of effector mechanism} & \textbf{Restriction factor} & \textbf{Restriction mechanism} & \textbf{Role as a disease driver} \\
    \hline
    \multirow{4}{2.6cm}{Nucleotide or nucleic acid degradation}
      & SAMHD1 \cite{cogginsSAMHD1FunctionsHuman2020,goldstoneHIV1RestrictionFactor2011,choiSAMHD1SpecificallyRestricts2015,orrisGuaninecontainingSsDNARNA2023,maharanaSAMHD1ControlsInnate2022,gutierrez-chamorroSAMHD1ExpressionModulates2023,felipModulationDNADamage2022} & Inhibition of cDNA synthesis, hydrolysis dNTPs & Aicardi--Goutières syndrome, depletion improves prognosis in some cancers \\
      \cline{2-4}
      & RNaseL \cite{klugeSnapShotAntiviralRestriction2015,maggHeterozygousOAS1Gainoffunction2021} & Degradation viral RNA & OAS1-mediated autoinflammatory immunodeficiency \\
      \cline{2-4}
      & ZAP/PARP13 \cite{shaoVersatilityZincFingerAntiviral2024} & Degradation viral RNA & Not reported \\
      \cline{2-4}
      & APOBEC family \cite{gottigApobec3ADeaminationFunctions2023,naumannDNADeaminationRequired2023} & Induction of hypermutation by deamination & Tumor development \\
    \hline
    \multirow{4}{2.6cm}{Transcription inhibition}
      & KAP1/TRIM28 \cite{xuIFI16PartnersKAP12022,leeCharacterizationInteractionTrim282018,shangSystematicPancancerAnalysis2023,whitworthWilmsTumourResulting2024,wuPhosphoproteomicsRevealsRole2022,yangKAP1StabilizesMYCN2024} & Gene silencing by recruitment of histone modifiers, DNA methylation & Higher tumor mutational burden and microsatellite instability, renal embryonal tumor, cell cycle progression in chronic lymphocytic leukemia, stabilization of MYCN mRNA \\
      \cline{2-4}
      & MxA \cite{turanNuclearMxAProteins2004,hayesExpressionMxAEsophageal2024,liExpressionMxAProtein2020} & GTPase activity, formation of a complex with viral nucleoprotein & Promotion of resistance to chemotherapy \\
      \cline{2-4}
      & MxB \cite{xieMxBImpedesNUP358mediated2020,wangInterferoninducibleMX2Host2020,schillingHumanMxBProtein2018} & GTPase activity & Not reported \\
      \cline{2-4}
      & IFI16 \cite{luIFI16PromotesProgression2024,xuIFI16PartnersKAP12022,howardNuclearDNASensor2022,johnsonIFI16RestrictsHSV12014} & Accumulation on the dsDNA viral genome & Promotion of clear renal cell carcinoma \\
    \hline
    \multirow{4}{2.6cm}{Translation inhibition}
      & SAMD9 \cite{mengParalogousPairMammalian2018,sivanHumanHostRange2018,suntharalinghamEmergingPhenotypesLinked2022} & Inhibition of ribosome entry and sequestration of viral RNA & MIRAGE syndrome \\
      \cline{2-4}
      & SAMD9L \cite{russellSAMD9LAutoinflammatoryAtaxia2021,mengParalogousPairMammalian2018,legrandSAMD9LActsAntiviral2024} & Restriction of viral protein synthesis through a Schlafen-like active site & Severe autoinflammatory disease and ataxia-pancytopenia \\
      \cline{2-4}
      & PKR \cite{morelPKRDoubleStranded2009,manjunathAPOBEC3BDrivesPKRmediated2023} & Phosphorylation of the eIF2$\alpha$ initiation factor & Deregulation of translation in Alzheimer's disease lymphocytes \\
      \cline{2-4}
      & SLFN11 \cite{valdezSchlafen11Restricts2019,okamotoSLFN11PromotesStalled2021a} & tRNA nucleolytic activity & Fanconi anemia \\
    \hline
    \multirow{4}{2.6cm}{Direct protein-virus interaction}
      & Tetherin/BST2/ CD317 \cite{ozcanBendingBST2Coiledcoil2017,tanwattanaHumanBST2Inhibits2023,jinBST2PromotesGrowth2019,heBST2InducedMacrophage2023} & Tethering of budding virus to the infected cell & Immunosuppressive tumor microenvironment, enhances cell proliferation and inhibits cell apoptosis in oral squamous cell carcinoma \\
      \cline{2-4}
      & IFITM family \cite{yuIFITMProteinsRestrict2015,hurInnateImmunityProtein2020} & Impairment of virion assembly through glycoprotein binding & Amyloid plaques formation \\
      \cline{2-4}
      & SERINC family \cite{leonhardtAntiviralHIV1SERINC2023} & Disruption of virus membrane asymmetry & Not reported \\
      \cline{2-4}
      & PSGL-1 \cite{liuPSGL1InhibitsHIV12020,wangRolePSGL1Pathogenesis2019} & Inhibition of viral assembly through sequestration of gp41 protein & Blockage increases survival in an endotoxemia mice model \\
    \hline
  \end{tabularx}
\end{table*}

%%%%%%%%%%%%%%%%%%%%%%%%%%%%%%%%%%%%%%%%%%%%%%%%%%%%%%%%%%%
\section{Concluding remarks}
%%%%%%%%%%%%%%%%%%%%%%%%%%%%%%%%%%%%%%%%%%%%%%%%%%%%%%%%%%%

Restriction mechanisms used by antiviral factors are able to target virtually all the steps of
the viral cycle and depend on physical interactions as well as enzymatic or signaling
activities (\crefcolor{tab:summary}). Because viral replication and host cellular systems
depend on the same molecular mechanisms, namely, replication, transcription and translation,
strategies used by cells to restrict viral infection can have deleterious effects on
physiological processes. However, restriction factors can also drive pathological phenotypes
through processes that are independent of their antiviral activity. In this review, we have
illustrated the main mechanisms by which most proteins of the autonomous innate immune system
can block the viral cycle. Additional mechanisms such as ISGylation of viral proteins or
modification of membrane fluidity have also been described as important levels of innate
cellular responses against certain viruses.

The restriction factors discussed in this review highlight the variety of biological
activities that can converge on the inhibition of a single molecular process. As our knowledge
of viral infectious mechanisms progresses, more cellular proteins that use novel mechanisms to
inhibit the viral cycle will be discovered. These studies will also evolve as we leverage new
technologies such as CRISPR-based screens that are already being used in the high-throughput
discovery of antiviral factors. Moreover, the disease-driving mechanisms discussed here stress
the importance of studying restriction factor function beyond immune response. An important
part of these proteins are involved in cellular processes outside of antiviral response and
can as such disturb key biological processes such as cell apoptosis and proliferation.
Overall, our understanding of restriction factors will be crucial to elucidate the complexity
of antiviral response and the role antiviral proteins and mechanisms can have in organisms as
a whole.

%%%%%%%%%%%%%%%%%%%%%%%%%%%%%%%%%%%%%%%%%%%%%%%%%%%%%%%%%%%
\section*{Acknowledgements}
%%%%%%%%%%%%%%%%%%%%%%%%%%%%%%%%%%%%%%%%%%%%%%%%%%%%%%%%%%%
I would like to thank François Roudier and Lucie Etienne for their valuable insight and advice
through the writing of this review. I would also like to thank Marine Tronchon for her helpful
comments on the manuscript.

%--- Bibliography ------------------------------------------------------------
\printbibliography[title=References]

% \balance equalises the two columns' height on this very last page — call
% it right after the bibliography prints, never earlier (it would disturb
% the layout of every page before it).
\balance

\end{document}
